\section{Methodik}
\label{chap:theorie}

Die DApp verlagert fachliche Kernregeln in Smart Contracts. Die Benutzeroberfläche
erlaubt Interaktionen. Die Zustandslogik liegt on-chain. Der Marktstatus ist
ständig aus der Blockchain ableitbar. Lesezugriffe erfolgen über RPC.
Schreibzugriffe erfolgen über signierte Transaktionen mit Gas-Kosten.

Für das Projekt bedeutet dies: Der Zustand bleibt transparent und nach
Deployment unveränderlich. Jede Änderung erfordert eine Wallet-Interaktion und
führt zu Transaktionskosten. Die Kopplung an Ereignisse oder externe
Entscheidungen ist modelliert, im vorliegenden System aber bewusst einfach.

Der ERC-20-Standard definiert eine Mindestschnittstelle für fungible Token:
Salden, Übertragungen und Genehmigungen. Das Projekt verwendet \texttt{VoteToken}
als einfaches Einsatzmittel. Der Token besitzt symbolischen Wert und verteilt sich
über ein einmaliges Faucet. So bleiben die relevanten Interaktionsschritte einer
typischen DApp sichtbar, ohne reales Vermögensrisiko.

Ein Predictionsmarkt sammelt Erwartungen über ein Ereignis in tokengewichteten
Pools. Teilnehmer ordnen Token einem Ausgang zu. Dadurch entsteht ein Verhältnis
zwischen den Pools. Dieses Verhältnis dient als Anzeige in der Benutzeroberfläche
und als Grundlage für die Abrechnung.

Märkte entstehen mit Frage, Kontext, Frist und Ausgängen. Teilnehmer staken Token
auf einen Ausgang. Nach Fristende setzt die Admin-Adresse den gewinnenden
Ausgang. Die Benutzeroberfläche reduziert neue Märkte auf die beiden Ausgänge
\emph{YES} und \emph{NO}, um den Ablauf kurz und nachvollziehbar zu halten.